% Thermodynamics
% Converted on 2026-10-08 from docs/lswt/02-observables/thermodynamics.md (status: draft, source-section: "Thermodynamics in Linear Spin Wave Theory: Partition Function, Internal Energy, Free Energy, Entropy Expression, Specific Heat (source pp. 13–15)").
% Review status: draft; awaiting the user's physics and mathematics review.

\hypertarget{sec-lswt-thermodynamics}{%
\section{Thermodynamics}\label{sec-lswt-thermodynamics}}

In thermal equilibrium the LSWT Hamiltonian describes an ideal gas of magnons on top of the classical reference configuration. The partition function therefore factorizes over the magnon modes \((\mathbf k,n)\), and the free energy, internal energy, entropy, and heat capacity follow from single-oscillator expressions summed over the magnetic Brillouin zone (MBZ). Energy, temperature, and entropy symbols follow \lswtnoteref{LSWT spin Hamiltonian}{Notation and Conventions}; the inverse temperature is \(\beta=(k_{\mathrm B}T)^{-1}\) and the entropy is \(\mathcal S\).

\subsection{Validity of the Magnon-Gas Description}

We use the diagonal Hamiltonian of \lswtnoteref{LSWT paraunitary diagonalization}{Paraunitary Diagonalization},

\[
\hat H_{\mathrm{LSWT}}
=E_{\mathrm{cl}}+\hat H_2
=E_{\mathrm{GS}}+\sum_{\mathbf k\in\mathrm{MBZ}}\sum_{n=1}^{N_{\mathrm{sub}}}\varepsilon_{n\mathbf k}\,\hat n_{\mathbf kn},
\qquad
\hat n_{\mathbf kn}=\hat b_{\mathbf kn}^\dagger\hat b_{\mathbf kn},
\]

with \(E_{\mathrm{GS}}=E_{\mathrm{cl}}+\Delta E_{\mathrm{zp}}\). The classical reference configuration and the magnon energies \(\varepsilon_{n\mathbf k}\) are held fixed as the temperature changes, and the quartic and higher Holstein--Primakoff terms that couple magnons are omitted. The results below are therefore the leading low-temperature description. Its corrections grow with the thermal magnon density, and the description fails when the thermal reduction of a sublattice moment, discussed in Magnon Observables (Section~\ref{sec-lswt-magnon-observables}), becomes comparable to the spin length. In two dimensions, the interpretation of these results as properties of an ordered phase requires the qualifications stated in \lswtnoteref{LSWT spin Hamiltonian}{LSWT Overview}.

All quantities below assume \(\varepsilon_{n\mathbf k}>0\), which holds when \(\mathsf H_{\mathbf k}\) is positive definite. The section on zero modes describes the limit \(\varepsilon_{n\mathbf k}\to0\).

\subsection{Partition Function}

The Gibbs state \(\hat\rho=\exp(-\beta\hat H_{\mathrm{LSWT}})/Z\) has the partition function \(Z=\operatorname{Tr}\exp(-\beta\hat H_{\mathrm{LSWT}})\). The number operators of different modes commute, so the trace factorizes into independent geometric series, one for each mode:

\begin{equation}\label{eq-lswt-partition-function}
Z=\exp(-\beta E_{\mathrm{GS}})\prod_{\mathbf k\in\mathrm{MBZ}}\prod_{n=1}^{N_{\mathrm{sub}}}\frac{1}{1-\exp(-\beta\varepsilon_{n\mathbf k})},
\qquad
-\ln Z=\beta E_{\mathrm{GS}}+\sum_{\mathbf k\in\mathrm{MBZ}}\sum_{n=1}^{N_{\mathrm{sub}}}\ln\!\left(1-\exp(-\beta\varepsilon_{n\mathbf k})\right).
\end{equation}

Each geometric series converges because \(\varepsilon_{n\mathbf k}>0\). The constant \(E_{\mathrm{GS}}\) contributes the factor \(\exp(-\beta E_{\mathrm{GS}})\); it shifts the internal and free energies but drops out of the entropy and heat capacity, which depend only on the thermal part of \(\ln Z\).

The thermal occupation of mode \((\mathbf k,n)\) is the Bose--Einstein distribution

\begin{equation}\label{eq-lswt-bose-einstein-distribution}
\langle\hat n_{\mathbf kn}\rangle
=n_{\mathrm B}(\varepsilon_{n\mathbf k}),
\qquad
n_{\mathrm B}(\varepsilon)=\frac{1}{\exp(\beta\varepsilon)-1}.
\end{equation}

\subsection{Internal and Free Energies}

The internal energy is the thermal average of the Hamiltonian, \(U=\langle\hat H_{\mathrm{LSWT}}\rangle=-\partial\ln Z/\partial\beta\). Differentiating \eqref{eq-lswt-partition-function} gives

\begin{equation}\label{eq-lswt-internal-energy}
U=E_{\mathrm{GS}}+\sum_{\mathbf k\in\mathrm{MBZ}}\sum_{n=1}^{N_{\mathrm{sub}}}\varepsilon_{n\mathbf k}\,n_{\mathrm B}(\varepsilon_{n\mathbf k}).
\end{equation}

The internal energy thus consists of the classical energy, the zero-point correction, and the energy of the thermally excited magnons. The Helmholtz free energy, defined by \(Z=\exp(-\beta F)\), is

\begin{equation}\label{eq-lswt-free-energy}
F=E_{\mathrm{GS}}+k_{\mathrm B}T\sum_{\mathbf k\in\mathrm{MBZ}}\sum_{n=1}^{N_{\mathrm{sub}}}\ln\!\left(1-\exp(-\beta\varepsilon_{n\mathbf k})\right).
\end{equation}

Each logarithm is negative, so the thermal magnons lower \(F\) below \(E_{\mathrm{GS}}\). At \(T=0\) both \(U\) and \(F\) reduce to \(E_{\mathrm{GS}}\).

\subsection{Entropy}

The entropy of the Gibbs state is the von Neumann entropy \(\mathcal S=-k_{\mathrm B}\operatorname{Tr}(\hat\rho\ln\hat\rho)\). Inserting \(\ln\hat\rho=-\beta\hat H_{\mathrm{LSWT}}-\ln Z\) gives \(\mathcal S=(U-F)/T\), which coincides with the thermodynamic relation \(\mathcal S=-\partial F/\partial T\). With \eqref{eq-lswt-internal-energy} and \eqref{eq-lswt-free-energy},

\[
\mathcal S=k_{\mathrm B}\sum_{\mathbf k\in\mathrm{MBZ}}\sum_{n=1}^{N_{\mathrm{sub}}}
\left[\beta\varepsilon_{n\mathbf k}\,n_{\mathrm B}(\varepsilon_{n\mathbf k})-\ln\!\left(1-\exp(-\beta\varepsilon_{n\mathbf k})\right)\right].
\]

The relations \(\beta\varepsilon=\ln[(1+n_{\mathrm B})/n_{\mathrm B}]\) and \(\ln(1-\exp(-\beta\varepsilon))=-\ln(1+n_{\mathrm B})\), which follow from the definition of \(n_{\mathrm B}\), express the entropy through the occupations alone:

\begin{equation}\label{eq-lswt-magnon-entropy}
\mathcal S=k_{\mathrm B}\sum_{\mathbf k\in\mathrm{MBZ}}\sum_{n=1}^{N_{\mathrm{sub}}}
\Big[\big(1+n_{\mathbf kn}\big)\ln\big(1+n_{\mathbf kn}\big)-n_{\mathbf kn}\ln n_{\mathbf kn}\Big],
\qquad
n_{\mathbf kn}=n_{\mathrm B}(\varepsilon_{n\mathbf k}).
\end{equation}

This is the entropy of an ideal Bose gas with the occupations \(n_{\mathbf kn}\). The zero-point correction does not enter, since it is a temperature-independent constant.

\subsection{Heat Capacity}

The heat capacity is the temperature derivative of the internal energy, \(C=\partial U/\partial T\). Only the occupations depend on temperature, and \(\partial n_{\mathrm B}(\varepsilon)/\partial T=(\varepsilon/k_{\mathrm B}T^2)\exp(\beta\varepsilon)/(\exp(\beta\varepsilon)-1)^2\), so

\begin{equation}\label{eq-lswt-magnon-heat-capacity}
C=k_{\mathrm B}\sum_{\mathbf k\in\mathrm{MBZ}}\sum_{n=1}^{N_{\mathrm{sub}}}
\frac{(\beta\varepsilon_{n\mathbf k})^2}{4\sinh^2(\beta\varepsilon_{n\mathbf k}/2)}.
\end{equation}

Each mode contributes the heat capacity of a harmonic oscillator. The contribution approaches \(k_{\mathrm B}\) for \(\varepsilon_{n\mathbf k}\ll k_{\mathrm B}T\) and is exponentially small for \(\varepsilon_{n\mathbf k}\gg k_{\mathrm B}T\), so at low temperature the heat capacity is dominated by the lowest magnon energies. The relation \(C=T\,\partial\mathcal S/\partial T\) gives the same expression.

\subsection{Normalization and Momentum Sums}

The quantities above are extensive: the sums run over the \(N_{\mathrm{uc}}\) momenta of the MBZ and the \(N_{\mathrm{sub}}\) bands. Values per magnetic site follow by dividing by \(N_{\mathrm{site}}=N_{\mathrm{uc}}N_{\mathrm{sub}}\); for example, \(U/N_{\mathrm{site}}=\mathcal E_{\mathrm{GS}}+N_{\mathrm{site}}^{-1}\sum_{\mathbf k,n}\varepsilon_{n\mathbf k}n_{\mathrm B}(\varepsilon_{n\mathbf k})\). In the thermodynamic limit the momentum sum becomes an integral over the MBZ, \(N_{\mathrm{uc}}^{-1}\sum_{\mathbf k}\to A_{\mathrm{MBZ}}^{-1}\int_{\mathrm{MBZ}}d^2\mathbf k\), where \(A_{\mathrm{MBZ}}\) is the area of the MBZ. A finite momentum mesh approximates this integral.

\subsection{Zero Modes}

When \(\mathsf H_{\mathbf k}\) is only positive semidefinite, some magnon energy vanishes, for example at a Goldstone mode. Near an isolated zero at \(\mathbf k_0\) we assume \(\varepsilon_{n\mathbf k}\propto|\mathbf k-\mathbf k_0|^a\) with \(a>0\); linear (\(a=1\)) and quadratic (\(a=2\)) dispersions are the common cases. The summands behave as follows as \(\varepsilon\to0\):

\begin{longtable}[]{@{}
  >{\raggedright\arraybackslash}p{(\columnwidth - 4\tabcolsep) * \real{0.3333}}
  >{\raggedright\arraybackslash}p{(\columnwidth - 4\tabcolsep) * \real{0.3333}}
  >{\raggedright\arraybackslash}p{(\columnwidth - 4\tabcolsep) * \real{0.3333}}@{}}
\toprule\noalign{}
\begin{minipage}[b]{\linewidth}\raggedright
Quantity
\end{minipage} & \begin{minipage}[b]{\linewidth}\raggedright
Summand for \(\varepsilon\ll k_{\mathrm B}T\)
\end{minipage} & \begin{minipage}[b]{\linewidth}\raggedright
Limit at \(\varepsilon=0\)
\end{minipage} \\
\midrule\noalign{}
\endhead
\bottomrule\noalign{}
\endlastfoot
\(U\) & \(\varepsilon\,n_{\mathrm B}(\varepsilon)\to k_{\mathrm B}T\) & finite \\
\(C\) & \(k_{\mathrm B}\) & finite \\
\(F\) & \(k_{\mathrm B}T\ln(\beta\varepsilon)\) & logarithmically divergent \\
\(\mathcal S\) & \(k_{\mathrm B}[1+\ln(k_{\mathrm B}T/\varepsilon)]\) & logarithmically divergent \\
\end{longtable}

The logarithmic divergences are integrable. In two dimensions \(\int d^2\mathbf q\,\ln|\mathbf q|\) converges near \(\mathbf q=0\) for any \(a>0\), so \(F\), \(U\), \(\mathcal S\), and \(C\) remain finite in the thermodynamic limit. A zero mode at an isolated momentum has vanishing weight in the MBZ integral. On a finite mesh that contains \(\mathbf k_0\), the treatment of that single point is a discretization choice, and the continuum values do not depend on it. Quantities weighted by the occupations without a compensating factor of \(\varepsilon\), such as the boson number and the sublattice moment, can diverge instead; they are treated in Magnon Observables (Section~\ref{sec-lswt-magnon-observables}).

% Added on 2026-10-10 (Claude, development decision D49). Review status: draft; awaiting the user's physics and mathematics review.
\subsection{Magnetization at Nonzero Temperature}\label{sec-lswt-thermal-magnetization}

A uniform Zeeman field of strength \(h\) along the unit vector \(\mathbf e\) enters the spin Hamiltonian as \(-h\sum_I\mathbf e\cdot\mathbf S_I\); a \(g\)-tensor replaces \(\mathbf e\) by \(\mathsf g^{\mathsf T}\mathbf e\) and changes nothing below. We define the magnetization per site at temperature \(T\) as the field derivative of the free energy \eqref{eq-lswt-free-energy},
\begin{equation}\label{eq-lswt-thermal-magnetization}
M(h,T)=-\frac{1}{N_{\mathrm{site}}}\,\frac{\mathrm dF}{\mathrm dh},
\end{equation}
in which the classical reference configuration at each \(h\) is the minimum of \(E_{\mathrm{cl}}\) at that field. At \(T=0\), \(F=E_{\mathrm{GS}}\) and \eqref{eq-lswt-thermal-magnetization} reduces to the zero-temperature magnetization \(-\mathrm d(E_{\mathrm{cl}}+\Delta E_{\mathrm{zp}})/\mathrm dh\) of Magnon Observables (Section~\ref{sec-lswt-magnon-observables}).

\subsubsection{Order of the Reference Configuration}

Holding the reference configuration at the minimum of \(E_{\mathrm{cl}}\) rather than of \(F\) is consistent to the order of LSWT. Write \(F(\boldsymbol\theta)=E_{\mathrm{cl}}(\boldsymbol\theta)+F_{\mathrm{qm}}(\boldsymbol\theta)\), where \(\boldsymbol\theta\) collects the classical angles and \(F_{\mathrm{qm}}=F-E_{\mathrm{cl}}\) is the zero-point and thermal magnon free energy of the quadratic Hamiltonian built at \(\boldsymbol\theta\). The classical energy is of order \(S^2\) and \(F_{\mathrm{qm}}\) of order \(S\) at fixed \(k_{\mathrm B}T/S\). Minimizing \(F\) instead of \(E_{\mathrm{cl}}\) shifts the angles by
\begin{equation}\label{eq-lswt-thermal-angle-shift}
\delta\boldsymbol\theta(T)=-\big(\partial_{\boldsymbol\theta}^2E_{\mathrm{cl}}\big)^{-1}\partial_{\boldsymbol\theta}F_{\mathrm{qm}},
\end{equation}
of order \(S^{-1}\), and lowers \(F\) by \(\tfrac12\,\partial_{\boldsymbol\theta}F_{\mathrm{qm}}\cdot(\partial_{\boldsymbol\theta}^2E_{\mathrm{cl}})^{-1}\partial_{\boldsymbol\theta}F_{\mathrm{qm}}\), of order \(S^0\). This is the order of the quartic Holstein--Primakoff terms that LSWT omits, so \eqref{eq-lswt-thermal-magnetization} is complete to order \(S^0\) without reminimizing the angles. Away from the classical minimum the boson Hamiltonian also has linear terms; they contribute to \(F\) only at second order in their coefficients, which vanish at the minimum, so they do not affect the first derivatives used below.

\subsubsection{Reduced Moment and Canting-Angle Shift}

The total derivative in \eqref{eq-lswt-thermal-magnetization} separates into three parts:
\[
\frac{\mathrm dF}{\mathrm dh}
=\frac{\partial E_{\mathrm{cl}}}{\partial h}\bigg|_{\boldsymbol\theta}
+\frac{\partial F_{\mathrm{qm}}}{\partial h}\bigg|_{\boldsymbol\theta}
+\partial_{\boldsymbol\theta}F_{\mathrm{qm}}\cdot\frac{\mathrm d\boldsymbol\theta_{\mathrm{cl}}}{\mathrm dh},
\]
since \(\partial_{\boldsymbol\theta}E_{\mathrm{cl}}=0\) at the classical minimum. The first term is \(-S_\mu\,\mathbf e\cdot\mathbf n_\mu\) summed over sites. For the second, the Hellmann--Feynman theorem for the Gibbs state gives \(\partial F_{\mathrm{qm}}/\partial h|_{\boldsymbol\theta}=\langle\partial\hat H_2/\partial h\rangle\). In the local frame \(\mathbf e\cdot\mathbf S_I=(\mathbf e\cdot\mathbf n_\mu)(S_\mu-\hat n_I)\) plus transverse components linear in the bosons, so the quadratic part of \(\partial\hat H/\partial h\) is \(\sum_I(\mathbf e\cdot\mathbf n_\mu)\hat n_I\) and the second term is \(\sum_I(\mathbf e\cdot\mathbf n_\mu)\langle\hat n_\mu\rangle_T\). For the third, differentiating \(\partial_{\boldsymbol\theta}E_{\mathrm{cl}}=0\) with respect to \(h\) gives \(\mathrm d\boldsymbol\theta_{\mathrm{cl}}/\mathrm dh=-(\partial_{\boldsymbol\theta}^2E_{\mathrm{cl}})^{-1}\partial_h\partial_{\boldsymbol\theta}E_{\mathrm{cl}}\) with \(\partial_h\partial_{\boldsymbol\theta}E_{\mathrm{cl}}=-\sum_IS_\mu\,\mathbf e\cdot\partial_{\boldsymbol\theta}\mathbf n_\mu\). With \eqref{eq-lswt-thermal-angle-shift} and the symmetry of \(\partial_{\boldsymbol\theta}^2E_{\mathrm{cl}}\),
\begin{equation}\label{eq-lswt-thermal-magnetization-decomposition}
M(h,T)=\frac{1}{N_{\mathrm{sub}}}\sum_\mu\mathbf e\cdot\Big[\big(S_\mu-\langle\hat n_\mu\rangle_T\big)\,\mathbf n_\mu+S_\mu\,\frac{\partial\mathbf n_\mu}{\partial\boldsymbol\theta}\cdot\delta\boldsymbol\theta(T)\Big].
\end{equation}
This is the zero-temperature moment of Magnon Observables (Section~\ref{sec-lswt-magnon-observables}) with the boson number and the angle shift replaced by their thermal values. For a collinear state along \(\mathbf e\) whose angles are fixed by symmetry, \(\partial_{\boldsymbol\theta}F_{\mathrm{qm}}=0\) and \(M\) is the thermal reduced moment alone. For a canted state the second term is of the same order \(S^0\) as the reduction \(\langle\hat n_\mu\rangle_T\), and the reduced moment alone does not give the magnetization at any temperature.

\subsubsection{Goldstone Modes in Two Dimensions}

When the field leaves a continuous symmetry about \(\mathbf e\), as for an isotropic Heisenberg model, the reference configuration breaks it and a Goldstone mode stays at zero energy for every \(h\). In two dimensions \(\langle\hat n_\mu\rangle_T\) then diverges logarithmically at \(T>0\) (Section~\ref{sec-lswt-magnon-observables}), and so does \(\partial_{\boldsymbol\theta}F_{\mathrm{qm}}\): a rotation of the angles away from the minimum gaps or destabilizes the Goldstone mode, whose thermal free energy depends on that gap as \(k_{\mathrm B}T\ln\varepsilon\). The two terms of \eqref{eq-lswt-thermal-magnetization-decomposition} diverge with opposite signs, and \(M\) is finite. The finiteness follows directly from \eqref{eq-lswt-thermal-magnetization}: the thermal part of \(\mathrm dF/\mathrm dh\) along the classical minimum is \(\sum_{\mathbf k,n}n_{\mathrm B}(\varepsilon_{n\mathbf k})\,\mathrm d\varepsilon_{n\mathbf k}/\mathrm dh\), and for a Goldstone mode \(\varepsilon_{n\mathbf k_0+\mathbf q}=c(h)|\mathbf q|\) the summand tends to the constant \(k_{\mathrm B}T\,c'(h)/c(h)\). The magnetization along the field is therefore a finite low-temperature quantity of LSWT in two dimensions, whereas the moment of each sublattice is not.

\subsubsection{Example: Canted Square-Lattice Antiferromagnet}

For the nearest-neighbour Heisenberg antiferromagnet on the square lattice with \(J=1\), lattice constant \(1\), spin \(S\), and \(\mathbf e=\hat{\mathbf z}\), the two sublattices cant toward the field by the angle \(\theta\) out of the plane perpendicular to \(\mathbf e\). In the frame rotated with each sublattice, one boson per site and momenta over the full square Brillouin zone describe the quadratic Hamiltonian at any \(\theta\):
\[
A_{\mathbf k}(\theta,h)=-4S\big(2\sin^2\theta-1\big)+h\sin\theta+4S\sin^2\theta\,\gamma_{\mathbf k},
\qquad
B_{\mathbf k}(\theta)=-4S\cos^2\theta\,\gamma_{\mathbf k},
\]
with \(\gamma_{\mathbf k}=\tfrac12(\cos k_x+\cos k_y)\), \(\varepsilon_{\mathbf k}=(A_{\mathbf k}^2-B_{\mathbf k}^2)^{1/2}\), and \(E_{\mathrm{cl}}/N_{\mathrm{site}}=2S^2(2\sin^2\theta-1)-hS\sin\theta\). The classical minimum \(\sin\theta=h/8S\) gives \(A_{\mathbf k}=4S(1+\sin^2\theta\,\gamma_{\mathbf k})\) and the classical magnetization \(h/8\) up to the saturation field \(h=8S\). The field enters \(A_{\mathbf k}\) only through \(h\sin\theta\), which makes the Hellmann--Feynman step of \eqref{eq-lswt-thermal-magnetization-decomposition} explicit: \(\partial F_{\mathrm{qm}}/\partial h|_\theta=N_{\mathrm{site}}\sin\theta\,\langle\hat n\rangle_T\) and \(\mathrm d\theta_{\mathrm{cl}}/\mathrm dh=(8S\cos\theta)^{-1}\).

The Goldstone mode sits at \((\pi,\pi)\). With \(\gamma\simeq-1+q^2/4\), \(A-B\simeq Sq^2\) and \(A+B\simeq8S\cos^2\theta\), so \(\varepsilon\simeq cq\) with
\[
c^2=8S^2\cos^2\theta=8S^2-\frac{h^2}{8}.
\]
The thermal summand of \(\langle\hat n\rangle_T\) near \((\pi,\pi)\) is \((A/\varepsilon)\,n_{\mathrm B}\simeq k_{\mathrm B}T/(2Sq^2)\), and its momentum integral gives \(\langle\hat n\rangle_T\simeq(k_{\mathrm B}T/4\pi S)\ln(1/q_{\min})\) for a smallest momentum \(q_{\min}\). The linear Goldstone mode gives \(F_{\mathrm{th}}/N_{\mathrm{site}}\simeq-\zeta(3)(k_{\mathrm B}T)^3/(2\pi c^2)\), and with \(\partial_hc^2=-h/4\),
\begin{equation}\label{eq-lswt-canted-square-thermal-magnetization}
M(h,T)-M(h,0)\simeq\frac{\zeta(3)\,h\,(k_{\mathrm B}T)^3}{8\pi c^4}
\end{equation}
at low temperature. The magnetization rises as \(T^3\), while each sublattice moment decreases without bound as \(q_{\min}\to0\). The magnon at \(\mathbf k=0\), the uniform precession about the field, has \(\varepsilon=8S\sin\theta=h\) and \(\mathrm d\varepsilon/\mathrm dh=1\); it lowers \(M\) by a term of order \(\exp(-h/k_{\mathrm B}T)\), which \eqref{eq-lswt-canted-square-thermal-magnetization} omits.

\subsection{References}

\subsubsection{Internal Documents}

\begin{itemize}
\tightlist
\item
  \lswtnoteref{LSWT spin Hamiltonian}{Notation and Conventions}: defines \(\beta\), \(\mathcal S\), \(\varepsilon_{n\mathbf k}\), \(E_{\mathrm{cl}}\), \(\Delta E_{\mathrm{zp}}\), \(E_{\mathrm{GS}}\), and the system-size symbols.
\item
  \lswtnoteref{LSWT paraunitary diagonalization}{Paraunitary Diagonalization}: supplies the diagonal magnon Hamiltonian, the zero-point correction, and the positive-definiteness condition.
\item
  \lswtnoteref{LSWT spin Hamiltonian}{LSWT Overview}: states the finite-temperature qualification for two-dimensional ordered states.
\item
  Magnon Observables (Section~\ref{sec-lswt-magnon-observables}): uses the Bose--Einstein occupations for boson numbers and sublattice moments, and supplies the zero-temperature moment with the canting-angle shift that Magnetization at Nonzero Temperature extends to \(T>0\).
\item
  Thermodynamic Derivations (Section~\ref{sec-lswt-thermodynamic-derivations}): is reserved for longer entropy and correlation-matrix derivations.
\end{itemize}

\subsubsection{External Sources}

\begin{itemize}
\tightlist
\item
  None.
\end{itemize}
